\subsubsection{Low and intermediate caps}
\label{sec:opt:low-upper}

The first algorithm is formalized in \cref{alg:raw}.

\begin{algorithm}[H]
\caption{\textsc{Raw}$(J,u)$}
\label{alg:raw}
\begin{algorithmic}
\Statex \textbf{Input:} jobs $J$ and a finite common upper limit $u$.
\State Fix an arbitrary order $j_1,\ldots,j_n$.
\State \textbf{For} $i=1,\ldots,n$, execute $j_i$ untested for $u$ time.
\end{algorithmic}
\end{algorithm}

For $u\le \uDet{2}$, use \textsc{Raw}.  It is optimal when $u\le\uDet{1}$;
when $u>1$, every effective length $\min\{u,1+p_j\}$ from
\cref{eq:model-effective-length} is at least one, so the ratio is at most
$u$.

For the next interval, let $R=\Rquad(u)$, $s=u-1$, and define
\begin{equation}
 K=s^2+s+1,
 \quad D=K+sR,
 \quad b=\frac{K-s^2R}{D},
 \quad \chi=\frac{1-b}{R},
 \quad \alpha=\frac KD.
 \label{eq:opt:ute-parameters}
\end{equation}
The defining equation for $R$ gives
\begin{equation}
 0\le b<\chi<\alpha<1,
 \qquad
 \chi=\frac{s(s+1)}D,
 \qquad
 b=(s+1)\alpha-s;
 \label{eq:opt:ute-order}
\end{equation}
For completeness, these inequalities follow without numerical root
isolation.  Since $T_s$ is increasing on positive arguments,
\begin{equation}
 T_s\!\left(\frac K{s^2}\right)
 =\frac K{s^3}\bigl(-s^3+s^2+2s+1\bigr)\ge0
 \qquad(s\le s_0),
\end{equation}
and hence $b=(K-s^2R)/D\ge0$, with equality only at $s=s_0$.  Moreover,
\begin{equation}
 \alpha-\chi=\frac1D,
 \qquad
 \chi-b=\frac{s^2R-1}{D}>0.
\end{equation}
For the last sign,
\begin{equation}
 9T_s(5/3)=f(s):=6s^3-12s^2-2s-3<0.
 \label{eq:opt:T-five-thirds}
\end{equation}
throughout the present interval.  Indeed, on $[4/5,1]$ the polynomial is
decreasing and $f(4/5)=-1151/125$; on $[1,s_0]$ it has one interior minimum,
while $f(1)=-11$ and $f(s_0)<f(13/6)=-95/36$.  Here
$s_0<13/6$ follows by substituting $13/6$ in
$s^3-(s+1)^2$.  Thus $R>5/3$, while $s>4/5$, and
$s^2R>16/15$.  Finally $K<D$ gives $\alpha<1$.  In particular, $b=0$
occurs only at $u=\uDet{3}$.
The remaining two elementary branches use the single procedure in
\cref{alg:ute}.

\begin{algorithm}[H]
\caption{\textsc{ForcedPrefixUTE}$(J,u,b)$}
\label{alg:ute}
\begin{algorithmic}
\Statex \textbf{Input:} jobs $J$, finite $u>1$, and $b\in[0,1]$.
\State Fix a test order $j_1,\ldots,j_n$; set
       $k\gets\lfloor bn\rfloor$, $\theta_u\gets\min\{1,u-1\}$,
       and $Q\gets\varnothing$.
\State \textbf{For} $i=1,\ldots,n$ \textbf{do}:
\State \quad Test $j_i$ for one unit and observe $p_{j_i}$.
\State \quad \textbf{If} $i\le k$ or $p_{j_i}\le\theta_u$, process $j_i$
       immediately.
\State \quad \textbf{Otherwise}, add $j_i$ to $Q$.
\State Process the jobs of $Q$ in nondecreasing revealed processing time.
\end{algorithmic}
\end{algorithm}

On $\uDet{2}\le u\le \uDet{3}$, use
\textsc{ForcedPrefixUTE} with $b$ from \cref{eq:opt:ute-parameters}.  Rounding
$bn$ changes the cost by $O_u(n)$.

Here is the complete binary calculation.  Adjacent exchanges place as many
long jobs as possible in the forced prefix and put all suffix-long jobs
before suffix zeros.  If $x$ is the total long mass, then the prefix-long
mass is $a=\min\{b,x\}$ and the deferred-long mass is
$d=(x-b)_+$.  After division by $n^2$, the leading excess
$G(x)=A-R O$ is
\begin{equation}
G(x)=
\begin{cases}
\displaystyle
 \frac{1-R}{2}+u\left(x-\frac{x^2}{2}\right)
 -\frac{R s}{2}x^2,
 &0\le x\le b,\\[3mm]
\displaystyle
 \frac{1-R}{2}
 +u\left(b-\frac{b^2}{2}\right)
 +(1-b)(x-b)+\frac{s}{2}(x-b)^2
 -\frac{R s}{2}x^2,
 &b\le x\le1.
\end{cases}
\label{eq:opt:binary-upper-gap}
\end{equation}
The second quadratic is strictly concave, has its stationary point at
$x=\alpha$, and satisfies $G(\alpha)=0$; after substituting
\cref{eq:opt:ute-parameters}, the latter identity is exactly
\cref{eq:opt:Rquad-equation}.  Since $b<\alpha$, its right derivative at $b$
is positive.  The left derivative satisfies
$G'_-(b)=G'_+(b)+s(1-b)>0$, and the first quadratic is concave, so it is
increasing on $[0,b]$.  Thus $G(x)\le0$ for every $x\in[0,1]$, including
the limiting case $b=0$.  This proves the binary upper bound without a
hidden local-versus-global step.

For $\uDet{2}\le u\le2$, arbitrary values reduce to binary values by
independent Bernoulli extremalization, in the following order.  First raise
\emph{every} suffix-deferred value to $u$.  This cannot lower ALG, while its
offline effective length was already $u$.  Only after this deterministic
step, replace a prefix value $p$ by $u$ with probability $p/u$ and by zero
otherwise, and replace a suffix-immediate value $p\le u-1$ by $u$ with
probability $p$ and by zero otherwise.  Notice that the $u$ realization
of a suffix-immediate job becomes deferred, so status is not preserved.
Nevertheless the online check is pairwise.  For jobs $i$ before $j$, the
leading pair contribution is
\begin{equation}
 g_{ij}=\begin{cases}
 1+p_i,&i\text{ is immediate},\\
 2+p_j,&i\text{ is deferred and }j\text{ is immediate},\\
 2+\min\{p_i,p_j\},&i,j\text{ are deferred}.
 \end{cases}
 \label{eq:opt:pair-contribution}
\end{equation}
A rounded prefix job stays immediate and preserves its expected value.  If
$i$ is suffix-immediate, its rounded contribution is one when it becomes
zero and at least two when it becomes $u$, hence its expectation is at least
$1+p_i$.  If $i$ is deferred and $j$ is suffix-immediate, then $i$ has been
raised to $u$ and the expected pair contribution is
$2+up_j\ge2+p_j$; a deferred--deferred pair can only increase.  Thus
expected ALG does not decrease, up to its $O_u(n)$ diagonal terms.

The expected effective lengths are,
respectively,
\begin{equation}
 1+\frac{u-1}{u}p\le \min\{u,1+p\},
 \qquad
 1+(u-1)p\le1+p,
\end{equation}
and $u$ for a deferred value.  Since the minimum is jointly concave,
$\EE\min(\lambda_i',\lambda_j')\le\min(\lambda_i,\lambda_j)$.  Thus some binary realization
has competitive excess at least that of the original instance.

It remains to justify arbitrary values for $2\le u\le \uDet{3}$.  The following
quadratic lemma records the entire endpoint calculation.

\begin{lemma}[UTE endpoint game]
\label{lem:opt:ute-endpoint-game}
For $2\le u\le \uDet{3}$, UTE with the parameters
\cref{eq:opt:ute-parameters} satisfies
\begin{equation}
 \ALG\le\Rquad(u)\OPT+O_u(n).
\end{equation}
\end{lemma}

\begin{proof}
Fix the symbolic prefix/immediate/deferred decisions.  Convexity in one
coordinate reduces a prefix job to $p\in\{0,u\}$, a suffix-immediate job to
$p\in\{0,1\}$, and a suffix-deferred job to $p\in\{1^+,u\}$.  For the last
claim, on $1<p<u-1$ the deferred--deferred term has coefficient
$1-R<0$ multiplying a concave minimum; on $[u-1,u]$, OPT is flat and ALG
is nondecreasing.  Adjacent exchanges give the canonical endpoint order.

Normalize counts by $n$.  Let $a$ be prefix-$u$ mass, $d$ suffix-$u$ mass,
$t$ suffix-$1^+$ mass, and $m$ suffix-immediate-one mass.  Their constraints
are
\begin{equation}
 0\le a\le b,
 \qquad d,t,m\ge0,
 \qquad d+t+m\le1-b.
\end{equation}
Pairwise accounting using \cref{eq:opt:pair-contribution} gives
\begin{align}
 A={}&\frac12+(s+1)\left(a-\frac{a^2}{2}\right)
 +(1-b)(d+t+m)-\frac{m^2}{2}+\frac{s}{2}d^2,
 \label{eq:opt:ute-A}\\
 O={}&\frac12+\frac{(a+d+t+m)^2}{2}
 +\frac{s-1}{2}(a+d)^2.
 \label{eq:opt:ute-O}
\end{align}
Moving $m$ to $t$ does not change $O$ and increases $A$, so $m=0$.  Put
$x=a+d$ and write $F(x)=A-R O$ after optimizing the other variables.
The derivative of the gap in $t$ is
$1-b-R(x+t)$, so its maximizing value is $(\chi-x)_+$.  At fixed
$x,t$, moving mass from $d$ to $a$ leaves $O$ unchanged and changes $A$ at
rate $s(1-x)+b-a\ge0$.  Hence
\begin{equation}
 t=(\chi-x)_+,
 \qquad
 a=\min\{b,x\},
 \qquad d=(x-b)_+.
 \label{eq:opt:ute-reduction}
\end{equation}

For $x\ge\chi$, the gap $A-R O$ is the binary concave quadratic already
discussed; it is tangent to zero at $x=\alpha$.  For $x<\chi$, put
\begin{equation}
 C_0=-\frac{R-1}{2}+\frac{(1-b)^2}{2R},
 \qquad
 \gamma=s-R(s-1).
\end{equation}
Substitution gives
\begin{equation}
 A-R O
 =C_0+a(s+b-sx)-\frac{a^2}{2}+\frac{\gamma x^2}{2}.
 \label{eq:opt:ute-positive-t}
\end{equation}
One has $\gamma>0$ on $1\le s\le s_0$.  Indeed,
$T_s(2)=s(s^2-s+1)>0$, so for $s\le2$ the equation $T_s(R)=0$ gives
$R<2$, and hence $\gamma>2-s\ge0$.  For
$2\le s\le s_0<13/6$,
\begin{equation}
 16T_s(7/4)=12s^3-20s^2+s-4>0;
\end{equation}
the polynomial is $14$ at $s=2$ and has positive derivative thereafter.
Thus $R<7/4$ and $\gamma>7/4-3s/4>1/8$.

On $0\le x\le b$, one has $a=x$, and
\begin{equation}
 A-R O=C_0+(s+b)x
 -\frac{s+1+R(s-1)}2x^2.
 \label{eq:opt:ute-small-x}
\end{equation}
Its derivative is decreasing, and the derivative at $b$ multiplied by $D$
is
\begin{equation}
 R\bigl(1+s^2+R s^2(s-1)\bigr)>0,
\end{equation}
so the maximum is at $b$.  On $b\le x\le\chi$,
\cref{eq:opt:ute-positive-t} is convex and its maximum is at $b$ or $\chi$.
The endpoint $\chi$ lies on the binary branch.  Finally,
\begin{equation}
 F(b)-F(\chi)
 =\frac{\chi-b}{2}\,[2sb-\gamma(b+\chi)],
\end{equation}
and direct use of \cref{eq:opt:ute-parameters} gives
\begin{equation}
 D[\gamma(b+\chi)-2sb]
 =-s+(1+s-s^3)R+s^2(s-1)R^2>0.
 \label{eq:opt:ute-small-certificate}
\end{equation}
For a rational sign certificate, $T_s$ is increasing on positive arguments
and
\begin{equation}
 9T_s(5/3)=6s^3-12s^2-2s-3<0,
\end{equation}
so $R>5/3$.  If the right-hand side of
\cref{eq:opt:ute-small-certificate} is denoted by $q_s(R)$, then
\begin{equation}
 3\,\partial_R q_s(5/3)=7s^3-10s^2+3s+3>0.
\end{equation}
The last polynomial is increasing for $s\ge1$, so $q_s$ is increasing for
$R\ge5/3$.  Finally,
\begin{equation}
 9q_s(5/3)=10s^3-25s^2+6s+15>0.
\end{equation}
For the last sign, write $s=3/2+z$; the expression becomes
$3/2-3z/2+20z^2+10z^3$.  It is at least $3/2$ for
$-1/2\le z\le0$, and for $z\ge0$ it is larger than
$3/2-9/320$.  These inequalities prove that every endpoint gap is
nonpositive.  Diagonals and prefix rounding contribute $O_u(n)$.
\end{proof}
At $u=\uDet{3}$, $b=0$.  The same zero-prefix algorithm remains exact through
$u=\uDet{4}$ and continues with a $\varphi$ guarantee afterwards.  We state
the continuation because it makes the three-algorithm lower envelope in
\cref{fig:intro-curves-det} literal rather than schematic.

\begin{lemma}[Zero-prefix extension]
\label{lem:opt:zero-prefix}
For every $u\ge \uDet{3}$,
\textnormal{\textsc{ForcedPrefixUTE}}$(J,u,0)$ from \cref{alg:ute} is
\begin{equation}
 \max\left\{\varphi,1+\frac1{\sqrt{u-1}}\right\}
 \label{eq:opt:zero-prefix-full-guarantee}
\end{equation}
competitive up to $O_u(n)$.  Thus its guarantee is
$1+1/\sqrt{u-1}$ on $[\uDet{3},\uDet{4}]$ and $\varphi$ on
$[\uDet{4},\infty)$.
\end{lemma}

\begin{proof}
The same coordinatewise reduction leaves capped-deferred mass $d$,
boundary-deferred mass $t$, immediate-one mass $m$, and zeros.  With
$s=u-1$, these masses satisfy $d,t,m\ge0$ and $d+t+m\le1$, and the leading
costs are
\begin{equation}
 A=\frac12+d+t+m-\frac{m^2}{2}+\frac{s d^2}{2},
 \qquad
 O=\frac12+\frac{(d+t+m)^2}{2}+\frac{(s-1)d^2}{2}.
\end{equation}
Move $m$ to $t$ and put $y=d+t$.  For fixed $y$, the ratio is
fractional-linear in $d^2$, so its maximum is on one of two faces:
\begin{equation}
 \frac{1+2y}{1+y^2}\le\varphi
 \quad(d=0),
 \qquad
 1+\frac{2y}{1+s y^2}\le1+\frac1{\sqrt s}
 \quad(d=y).
\end{equation}
The first one-variable function is maximized at $y=1/\varphi$ and the second
at $y=1/\sqrt s$, which gives the two displayed constants.
The second value dominates the first exactly when
$u\le\uDet{4}$.  Binary instances attain the second value.
\end{proof}


\subsubsection{High caps}
\label{sec:optional-upper}

We now prove the upper bounds on the two ranges on which capped outcomes
matter.  Besides completing the upper-bound side of the revealing-optimization curve,
the proof records explicitly why an outcome at the raw cap cannot be
treated as an ordinary deferred outcome.  Throughout this subsection,
$R=1+c$, and we retain $\varphi$, $\uDet{5}$, and $r$ from
\cref{eq:opt:u1-u2,eq:opt:obligatory-constant}.

For every $c>0$ define the history-dependent threshold
\begin{equation}
 a_c(y)=
 \begin{cases}
  1, & y\le-1,\\[1mm]
  \displaystyle
  1+\frac12\ln\!\left(\frac{c+2}{c-2y}\right),
      & -1\le y\le0.
 \end{cases}
 \label{eq:opt:threshold}
\end{equation}
The single algorithm used throughout the last two branches is the following.
It does not use raw execution and therefore does not depend on $u$.

\begin{algorithm}[H]
\caption{\textsc{AdaptiveThreshold}$(J,c)$}
\label{alg:adaptive-threshold}
\begin{algorithmic}
\Statex \textbf{Input:} jobs $J$ and parameter $c>0$.
\State Fix an arbitrary test order $j_1,\ldots,j_n$.
\State Set $N\gets0$, $N_{\mathrm{def}}\gets0$, and $Q\gets\varnothing$.
\State \textbf{For} $i=1,\ldots,n$ \textbf{do}:
\State \quad Set $x\gets n-i+1$ and $y\gets(N_{\mathrm{def}}-(1+c)N)/x$.
\State \quad Set $a\gets a_c(y)$ according to \cref{eq:opt:threshold}.
\State \quad Test $j_i$ for one unit and observe $p_{j_i}$.
\State \quad \textbf{If} $p_{j_i}\le a$, process $j_i$ immediately.
\State \quad \textbf{Otherwise}, add $j_i$ to $Q$ and set $N_{\mathrm{def}}\gets N_{\mathrm{def}}+1$.
\State \quad \textbf{If} $p_{j_i}>0$, set $N\gets N+1$.
\State Process the jobs of $Q$ in nondecreasing order of their revealed
       processing times.
\end{algorithmic}
\end{algorithm}

For $c\in[r,1/\varphi]$, let $m=m(c)\in[0,1/\varphi]$ and $u=u(c)$ be
determined by the common parametrization
\cref{eq:opt:mixed-parametrization}.
The right-hand side of the first equation is strictly increasing in $c$,
as is the left-hand side in $m$.  Consequently $m(c)$ is increasing and
$u(c)$ is decreasing.  The endpoints are
\[
 (c,m,u)=\left(\frac1\varphi,\frac1\varphi,\uDet{4}\right)
 \quad\text{and}\quad
 (c,m,u)=(r,0,\uDet{5}).
\]

\begin{theorem}[Middle and high revealing upper bounds]
\label{thm:opt:optional-middle-upper}
For every finite $u\in[\uDet{4},\infty)$, the formally specified algorithm
\textnormal{\textsc{AdaptiveThreshold}} from
\cref{alg:adaptive-threshold}
satisfies
\[
 \ALG\le
 \begin{cases}
  (1+c)\OPT+O_u(n),
     & \uDet{4}\le u\le \uDet{5},
       \quad c\text{ given by \eqref{eq:opt:mixed-parametrization}},\\[1mm]
  (1+r)\OPT+O(n),
     & u\ge \uDet{5}.
 \end{cases}
\]
The constant in the last $O(n)$ is absolute and uniform over
$u\in[\uDet{5},\infty)$.
In particular these are upper bounds on the deterministic
size-asymptotic competitive ratio.
\end{theorem}

\begin{proof}

Fix a finite cap $u$ in the stated range.

\paragraph{The algorithm.}
Use \textsc{AdaptiveThreshold}$(J,c)$ from
\cref{alg:adaptive-threshold}; on the high interval use $c=r$.
For the analysis, immediately before a test let $x$ be the number of jobs not
yet tested, $N$ the number of earlier strictly positive outcomes, and $N_{\mathrm{def}}$ the
number of those outcomes deferred to the final tail.  Thus
\begin{equation}
 y=\frac{N_{\mathrm{def}}-RN}{x}\le0
 \label{eq:opt:state-y}
\end{equation}
and the threshold is given by \cref{eq:opt:threshold}.  On the mixed interval
$c\ge r$, and hence
\[
 a_c(y)\le a_c(0)\le a_r(0)=\frac1r<u-1.
 \label{eq:opt:threshold-below-cap}
\]
The same inequality holds on the high interval, where we use $c=r$.
Thus every outcome in the capped region $[u-1,u]$ is deferred.

The state $y$ never increases.  Indeed, the numerator in
\eqref{eq:opt:state-y} changes by $0$, $-R$, or $-c$ after a zero, a
positive immediate outcome, or a positive deferred outcome, respectively,
while $x$ decreases.  This monotonicity will allow us to allocate every
pair cost when its relevant endpoint becomes known.

\paragraph{Exact pair objective and the five endpoints.}
Write
\[
 \bar p_i=\min\{p_i,u-1\}.
\]
By \cref{eq:model-effective-length,lem:offline}, the clairvoyant revealing
optimum is SPT on jobs of effective lengths
$1+\bar p_i=\min\{1+p_i,u\}$.  For two jobs tested in the order $i<j$, let
\begin{equation}
 e_{ij}=
 \begin{cases}
  (p_i-p_j)_+,
     &i\text{ is processed immediately},\\
  1+(p_j-p_i)_+,
     &i\text{ is deferred and }j\text{ is immediate},\\
  1,
     &i\text{ and }j\text{ are deferred}.
 \end{cases}
 \label{eq:opt:pair-excess}
\end{equation}
The first line applies whenever $i$ is immediate, regardless of the
action taken on $j$.  Pairwise completion-time accounting for the online
schedule gives
\begin{equation}
 \ALG
 =\sum_i(1+p_i)
   +\sum_{i<j}\bigl(1+\min\{p_i,p_j\}+e_{ij}\bigr),
 \label{eq:opt:alg-pair-decomposition}
\end{equation}
whereas the common offline pair identity \cref{eq:optional-opt-pair} gives
\begin{equation}
 \OPT
 =\sum_i(1+\bar p_i)
   +\sum_{i<j}\bigl(1+\min\{\bar p_i,\bar p_j\}\bigr).
 \label{eq:opt:opt-pair-decomposition}
\end{equation}
It is therefore natural to introduce
\begin{equation}
 \Gamma_n=\sum_{i<j}
 \left(e_{ij}+\min\{p_i,p_j\}
       -R\min\{\bar p_i,\bar p_j\}\right).
 \label{eq:opt:pair-objective}
\end{equation}
Equations \eqref{eq:opt:alg-pair-decomposition}--
\eqref{eq:opt:pair-objective} yield the exact identity
\begin{equation}
 \ALG-R\OPT
 =\Gamma_n-c\binom n2
  +\sum_i\bigl[(1+p_i)-R(1+\bar p_i)\bigr].
 \label{eq:opt:gap-decomposition}
\end{equation}
The final sum is nonpositive.  Below the cap its $i$-th summand is
$-c(1+p_i)$; on the cap it is at most $1-cu\le0$.  Thus it is enough
to control $\Gamma_n$.

Fix a symbolic immediate/deferred path.  Coordinatewise convexity in
\eqref{eq:opt:pair-objective} reduces a worst outcome vector to the five
formal endpoint types
\begin{equation}
\begin{array}{c|c|c}
 \text{type}&\text{outcome}&\text{action}\\ \hline
 Z&0&\text{immediate}\\
 E&0^+&\text{immediate}\\
 M&a_c(y)&\text{immediate}\\
 Q&a_c(y)^+&\text{deferred}\\
 U&u&\text{deferred and capped}.
\end{array}
 \label{eq:opt:five-endpoints}
\end{equation}
For an immediate coordinate the two positive endpoints are $0^+$ and
the live threshold; the genuine zero is kept separate because it does not
enter $N$.  A deferred coordinate below $u-1$ moves to the live
threshold or to $u-1$.  On $[u-1,u]$, the offline contribution is flat
and the algorithmic contribution is nondecreasing, so the latter endpoint
moves to $u$.  The symbolic categories, and therefore all later
threshold states, remain unchanged during each coordinate move.  This is
the only endpoint reduction used below.

\paragraph{A local pair allocation.}
Separate the earlier $E$-outcomes from the substantive outcomes
$M,Q,U$.  Let
\[
 P=\#\{\text{earlier }M,Q,U\},\qquad
 E=\#\{\text{earlier }E\},\qquad
 N_{\mathrm{def}}=\#\{\text{earlier }Q,U\},
\]
and let $K$ be the number of earlier $U$-outcomes.  Among the ordinary
substantive outcomes, let $I,d$ be the numbers of $M,Q$-outcomes and
let $A_I,A_D$ be the sums of their processing values.  Thus
$P=I+d+K$ and $N_{\mathrm{def}}=d+K$.

At a current ordinary endpoint, with $a=a_c(y)$, allocate the relaxed
local charges
\begin{equation}
 \ell_Z=\ell_E=N_{\mathrm{def}},
 \qquad
 \ell_M=N_{\mathrm{def}}+a(x+N_{\mathrm{def}}-RP),
 \qquad
 \ell_Q=N_{\mathrm{def}}+a(N_{\mathrm{def}}-RP).
 \label{eq:opt:ordinary-local-charges}
\end{equation}
For a capped endpoint the exact charge is instead
\begin{align}
 \ell_U
 &=N_{\mathrm{def}}+A_D+uK-R\bigl(A_I+A_D+(u-1)K\bigr)
 \notag\\
 &=d-RA_I-cA_D+K h_U,
 \qquad
 h_U=1+u-R(u-1).
 \label{eq:opt:cap-local-charge}
\end{align}
These formulae follow by checking the finitely many ordered endpoint
pairs.  In particular, an earlier immediate outcome already allocated its
online processing delay when it appeared, so a current cap contributes
only $-Rp$ against that outcome.  Against an earlier ordinary deferred
outcome it contributes $p-Rp=-cp$, and against an earlier cap it
contributes $h_U$.  The same check shows
\begin{equation}
 \Gamma_n\le\sum_i\ell_i;
 \label{eq:opt:pair-allocation-dominates}
\end{equation}
the difference is exactly the sum of the nonnegative own-processing terms
allocated by the $M$-endpoints.  Formula
\eqref{eq:opt:cap-local-charge}, rather than a substitution $p=u$ into
$\ell_Q$, is essential before caps have been put into a prefix.

Threshold monotonicity gives $A_I\ge aI$ and $A_D\ge ad$.  Comparing
\eqref{eq:opt:cap-local-charge} with the hypothetical $Q$-charge at the
same state gives the exact identity
\begin{align}
 \ell_U-\ell_Q
 &=R(aI-A_I)+c(ad-A_D)
   +K\,[1-c(u-1-a)]
 \notag\\
 &\le K\,[1-c(u-1-a)].
 \label{eq:opt:cap-versus-q}
\end{align}

\paragraph{The reusable four-endpoint bank.}
The states $Z,E,M,Q$ are the four ordinary, noncap endpoints.  The potential
below proves their inequalities simultaneously for every target parameter
$c$ on the mixed interval.  The fifth endpoint $U$ is handled only through
the extra reserve.  At $c=r$ that reserve vanishes, leaving a cap-free bank
that will also support the obligatory endpoint without a second proof.
Define
\begin{equation}
 \eta=\frac{N_{\mathrm{def}}-RP}{x},
 \qquad b=\frac Ex,
 \qquad y=\eta-Rb.
 \label{eq:opt:eta-b}
\end{equation}
For $-1\le y\le0$, put
\begin{align}
 f_c(y)
 &=\frac{c+2}{4}
   +\frac{c-2y}{4}
      \left(\ln\frac{c-2y}{c+2}-1\right),
 \label{eq:opt:f-definition}\\
 h_c(y)&=a_c(y)(1+y),
 \notag\\
 G_c(y,b)&=f_c(y)+b h_c(y)+\frac R2b^2.
 \label{eq:opt:G-region-one}
\end{align}
For $y\le-1$, set instead
\begin{equation}
 G_c(y,b)=g_c(\eta),
 \qquad
 g_c(\eta)=\frac{(1+\eta)_+^2}{2R}.
 \label{eq:opt:G-region-two}
\end{equation}
The two definitions are $C^1$ across $y=-1$.  Indeed, feasibility
implies $0\le b\le1/R$ there, and both sides have value $Rb^2/2$ and
derivatives
\[
 (G_y,G_b)=(b,Rb).
\]
They are also nonnegative: $f_c(-1)=0$,
$f_c'=a_c-1\ge0$, and $h_c,b,g_c\ge0$.

Use the base potential
\begin{equation}
 W_{\mathrm{base}}(x,P,E,N_{\mathrm{def}})=xN_{\mathrm{def}}+x^2G_c(y,b).
 \label{eq:opt:base-potential}
\end{equation}
In the raw state $(x,P,E,N_{\mathrm{def}},K)$, the endpoint directions are
\begin{equation}
\begin{array}{c|ccccc}
 &x&P&E&N_{\mathrm{def}}&K\\ \hline
 v_Z&-1&0&0&0&0\\
 v_E&-1&0&1&0&0\\
 v_M&-1&1&0&0&0\\
 v_Q&-1&1&0&1&0\\
 v_U&-1&1&0&1&1.
\end{array}
 \label{eq:opt:state-directions}
\end{equation}
Let $J_T=\ell_T+\nabla W_{\mathrm{base}}\mathbin{\cdot}v_T$, where for
the moment $T\in\{Z,E,M,Q\}$.  In Region I, direct differentiation and
the identities
\begin{align}
 -2f_c+(y-R)f_c'+a_c(1+y)&=0,
 \notag\\
 1-2f_c+(y-c)f_c'+a_cy&=0
 \label{eq:opt:f-hjb-identities}
\end{align}
give
\begin{align}
 \frac{J_Z}{x}
 &=(-2f_c+yf_c')+b(-h_c+yh_c'),
 \notag\\
 \frac{J_E}{x}
 &=b[-h_c+(y-R)h_c'+R],
 \notag\\
 \frac{J_M}{x}
 &=b[-h_c+(y-R)h_c'+a_cR],
 \notag\\
 \frac{J_Q}{x}
 &=b[-h_c+(y-c)h_c'+a_cR].
 \label{eq:opt:region-one-drifts}
\end{align}
All four quantities are nonpositive.  For the first one,
\begin{equation}
 -2f_c+yf_c'=a_c(c-y)-R\le0;
 \label{eq:opt:zero-drift-sign}
\end{equation}
the right-hand side is zero at $y=-1$ and is nonincreasing because
$a_c'(c-y)-a_c\le1-a_c\le0$.  Also $h_c,h_c'\ge0$.  For the remaining
three drifts, observe that
\begin{equation}
 h_c+(c-y)h_c'
 =a_cR+(c-y)(1+y)a_c'\ge a_cR.
 \label{eq:opt:h-sign-identity}
\end{equation}
The $Q$-bracket in \eqref{eq:opt:region-one-drifts} is therefore
nonpositive; the $M$-bracket is the $Q$-bracket minus $h_c'$, and the
$E$-bracket is the $M$-bracket minus $R(a_c-1)$.

In Region II, the same calculation gives
\begin{align}
 J_Z/x=J_E/x&=-2g_c+\eta g_c',
 \notag\\
 J_M/x&=-2g_c+(\eta-R)g_c'+1+\eta,
 \notag\\
 J_Q/x&=1-2g_c+(\eta-c)g_c'+\eta.
 \label{eq:opt:region-two-drifts}
\end{align}
For $-1\le\eta\le0$, substitution of
$g_c'=(1+\eta)/R$ makes these values $-(1+\eta)/R$,
$-(1+\eta)/R$, and $0$, respectively.  For $\eta\le-1$ their signs
are immediate from $g_c=g_c'=0$.  Hence the base bank pays every
ordinary endpoint.

\paragraph{The cap reserve.}
Define
\begin{equation}
 Z(c)=2+\frac1c+\frac12\ln\!\left(1+\frac2c\right),
 \qquad
 \delta=Z(c)-u.
 \label{eq:opt:Z-delta}
\end{equation}
On the mixed curve, \eqref{eq:opt:mixed-parametrization} says precisely
that
\begin{equation}
 \delta=\operatorname{artanh}m.
 \label{eq:opt:delta-atanh}
\end{equation}
For $0\le t\le m$, set
\begin{equation}
 \beta(t)=ct\,[\delta-\operatorname{artanh}t],
 \qquad
 \mathcal V(t)=\int_t^m\beta(s)\,ds,
 \label{eq:opt:reserve-integral}
\end{equation}
and extend $\mathcal V$ by zero for $t\ge m$.  Since
$\beta(m)=0$, this extension is $C^1$.  With
\begin{equation}
 q=\frac Kx,
 \qquad
 \mu=\frac{q}{1+q},
 \qquad
 \mathcal H(q)=(1+q)^2\mathcal V(\mu),
 \label{eq:opt:reserve-H}
\end{equation}
define the full potential
\begin{equation}
 W=xN_{\mathrm{def}}+x^2\bigl(G_c(y,b)+\mathcal H(q)\bigr).
 \label{eq:opt:full-potential}
\end{equation}
It is nonnegative.

For an ordinary endpoint, $K$ is fixed while $x$ decreases, so the
reserve contributes, after division by $x$,
\begin{equation}
 L_0(q)=-2\mathcal H(q)+q\mathcal H'(q).
 \label{eq:opt:L0}
\end{equation}
A cap also increments $K$, and its contribution is
$L_0+\mathcal H'$.  On the active range $\mu\le m$, differentiation of
\eqref{eq:opt:reserve-H} gives
\begin{equation}
 \mathcal H'=2(1+q)\mathcal V-\beta(\mu),
 \qquad
 L_0=-2(1+q)\mathcal V-q\beta(\mu)\le0.
 \label{eq:opt:reserve-derivatives}
\end{equation}
Thus the reserve can only improve every ordinary drift.

For a cap, feasibility $N_{\mathrm{def}}\le P$ and $P\ge K$ gives
\begin{equation}
 y=\frac{N_{\mathrm{def}}-R(P+E)}x\le-\frac{cK}{x}=-cq.
 \label{eq:opt:y-cap-bound}
\end{equation}
While $\mu\le m$, we have $q\le m/(1-m)\le\varphi$ and hence
$cq\le1$.  Monotonicity of $a_c$, together with
\eqref{eq:opt:cap-versus-q}, yields
\begin{align}
 \frac{\ell_U-\ell_Q}{x}
 &\le cq\left[\delta-\frac12\ln(1+2q)\right]
 =:B(q).
 \label{eq:opt:cap-residual-B}
\end{align}
The change of variables in \eqref{eq:opt:reserve-H} satisfies
\begin{equation}
 \operatorname{artanh}\mu=\frac12\ln(1+2q),
 \qquad
 (1+q)\beta(\mu)=B(q).
 \label{eq:opt:B-beta}
\end{equation}
Consequently the cap residual is cancelled exactly:
\begin{equation}
 L_0+\mathcal H'+B(q)=0
 \qquad(\mu\le m).
 \label{eq:opt:reserve-cancellation}
\end{equation}
When $\mu\ge m$, the reserve vanishes.  At
$q_*=m/(1-m)$, the bracket in
\eqref{eq:opt:cap-versus-q}, evaluated at the largest possible threshold
$a_c(-cq_*)$, is zero by
\eqref{eq:opt:delta-atanh} and \eqref{eq:opt:B-beta}.  For $q\ge q_*$,
\eqref{eq:opt:y-cap-bound} only decreases the threshold, so that bracket is
nonpositive.  The fifth Hamiltonian inequality therefore holds for every
state and for arbitrary interleavings of capped and ordinary outcomes.

\paragraph{Initial calibration.}
Initially $y=b=q=0$, and
\begin{equation}
 f_c(0)=\frac12-\frac c4\ln\!\left(1+\frac2c\right).
 \label{eq:opt:f-zero}
\end{equation}
Elementary integration in \eqref{eq:opt:reserve-integral} gives
\begin{equation}
 \mathcal V(0)=\frac c2(\delta-m).
 \label{eq:opt:V-zero}
\end{equation}
Indeed,

\[
 \int_0^m t\operatorname{artanh}t\,dt
 =\frac{(m^2-1)\operatorname{artanh}m+m}{2}.
\]
Using \eqref{eq:opt:mixed-parametrization} and
\eqref{eq:opt:delta-atanh}, we obtain the exact calibration
\begin{equation}
 f_c(0)+\mathcal H(0)
 =f_c(0)+\mathcal V(0)=\frac c2.
 \label{eq:opt:initial-calibration}
\end{equation}
Thus $W_{\mathrm{initial}}=cn^2/2$.

\paragraph{Finite jobs, and uniformity on the plateau.}
We spell out the regularity needed to pass from the fluid Hamiltonian to
unit job transitions.  It is important here to restrict attention to
feasible states.  In Region I, $y\ge-1$ implies
\begin{equation}
 -xy=R(P+E)-N_{\mathrm{def}}=cP+(P-N_{\mathrm{def}})+RE\le x.
 \label{eq:opt:feasible-compactness}
\end{equation}
All terms in the middle expression are nonnegative.  Hence
$P/x,N_{\mathrm{def}}/x,q\le1/c$ and $b\le1/R$, so every normalized derivative of
the Region-I perspective in \eqref{eq:opt:base-potential} is bounded on
the feasible cone.  In Region II the raw base bank is simply
\begin{equation}
 xN_{\mathrm{def}}+\frac{(x+N_{\mathrm{def}}-RP)_+^2}{2R},
 \label{eq:opt:region-two-raw-bank}
\end{equation}
which has bounded one-sided Hessians.  The two pieces are $C^1$ at
$y=-1$, as verified above, and the positive-part square is $C^1$ at
$\eta=-1$.

For the reserve $\mathcal R(x,K)=x^2\mathcal H(K/x)$, on its active cone
$0\le q\le q_*$ one has bounded $\mathcal H''$, and
\begin{equation}
 \mathcal R_{xx}=2\mathcal H-2q\mathcal H'+q^2\mathcal H'',
 \quad
 \mathcal R_{xK}=\mathcal H'-q\mathcal H'',
 \quad
 \mathcal R_{KK}=\mathcal H''.
 \label{eq:opt:reserve-hessian}
\end{equation}
For $q\ge q_*$, $\mathcal R=0$; moreover
$\mathcal H(q_*)=\mathcal H'(q_*)=0$.  Thus the join is $C^1$ with
bounded one-sided Hessians.  At $x=0$, set $\mathcal R=0$: for $K>0$ it is
already locally zero, while at $(x,K)=(0,0)$ two-homogeneity gives zero
gradient.  Altogether $W$ is $C^{1,1}$ on the feasible state cone.  On the
mixed interval its regularity constant may depend on the fixed value of $u$.
When $c=r$ and $m=0$, however, the reserve is identically zero and the base
bank depends only on the absolute constant $r$; its regularity constant is
therefore uniform over all $u\ge \uDet{5}$.  This includes the last transition
from $x=1$ to $x=0$.

It follows that a unit transition has the Taylor remainder
\begin{equation}
 \ell_i+W_{i+1}-W_i\le O_u(1).
 \label{eq:opt:finite-taylor}
\end{equation}
The endpoint $E=0^+$ is interpreted one-sidedly.  For a fixed symbolic
word, replace all $E$-coordinates by positive values $\eps$.  Its
state trajectory is unchanged, and the exact finite pair objective, being
a finite sum of affine and minimum terms, converges as
$\eps\downarrow0$.  Equivalently, retain an
$O(\eps n^2)$ error and take this limit for each fixed $n$.
Therefore the formal $E$-Hamiltonian transfers to genuine outcome
vectors without an asymptotic loss.

Summing \eqref{eq:opt:finite-taylor}, using $W\ge0$, and applying
\eqref{eq:opt:initial-calibration} gives
\begin{equation}
 \sum_i\ell_i\le\frac c2n^2+O_u(n).
 \label{eq:opt:ell-telescope}
\end{equation}
Together with \eqref{eq:opt:gap-decomposition} and
\eqref{eq:opt:pair-allocation-dominates}, this yields
\begin{align}
 \ALG-R\OPT
 &\le -c\binom n2+\frac c2n^2+O_u(n)
 =O_u(n).
 \label{eq:opt:mixed-upper-finite}
\end{align}

At $u=\uDet{5}$, one has $c=r$, $m=0$, and
$Z(r)=\uDet{5}$.  The cap reserve is identically zero, while
\eqref{eq:opt:initial-calibration} remains valid.  If $u$ is increased
above $\uDet{5}$, the only adverse fifth-endpoint bracket
$1-r(u-1-a_r(y))$ decreases, and the capped diagonal term in
\eqref{eq:opt:gap-decomposition} also decreases.

Here is the promised uniform finite calculation.  The base bank and its
feasible cone contain no $u$, so there is one constant $C_{\rm step}$, depending
only on $r$, for which every ordinary or capped transition with $u\ge \uDet{5}$
satisfies
\[
 \ell_i+W_{i+1}-W_i\le C_{\rm step}.
\]
Because $W_0=rn^2/2$ and $W_n=0$, telescoping gives
\begin{equation}
 \sum_i\ell_i\le\frac r2n^2+C_{\rm step}n.
 \label{eq:opt:uniform-ell}
\end{equation}
In the diagonal term of \cref{eq:opt:gap-decomposition}, a value below the
cap contributes $-r(1+p_i)$.  A capped value contributes at most
$1-r u\le0$.  Hence the diagonal sum is nonpositive uniformly on
$[\uDet{5},\infty)$.  Combining
\cref{eq:opt:gap-decomposition,eq:opt:pair-allocation-dominates,eq:opt:uniform-ell}
now yields
\begin{align}
 \ALG-(1+r)\OPT
 &\le-r\binom n2+\frac r2n^2+C_{\rm step}n
   =\left(C_{\rm step}+\frac r2\right)n.
 \label{eq:opt:uniform-high-gap}
\end{align}
Thus, with $B=C_{\rm step}+r/2$,
\begin{equation}
 \ALG\le(1+r)\OPT+B n
 \qquad\text{for every finite }u\in[\uDet{5},\infty),
 \label{eq:opt:uniform-high-upper}
\end{equation}
and $B$ is independent of $u$.  The estimate is genuinely uniform rather
than the result of sending a parameter-dependent remainder to infinity.  This proves
\cref{thm:opt:optional-middle-upper}.  Finally, because $u>1$ on the
present range, every effective offline length is at least one and
$\OPT\ge n(n+1)/2$.  The additive term in
\eqref{eq:opt:mixed-upper-finite} vanishes in the size-asymptotic ratio.

\end{proof}
